\chapter{Conclusion}
\label{chapterlabel6}

This dissertation examined how night-time public transport use in London is organised in time and space, how Rail and Bus differ in that organisation, and how the resulting usage types relate to night-work geography and wider socio-spatial context. Separate mode-specific typologies were built for the 18:00--05:00 window. The findings show that night-time public transport use is a temporally differentiated system rather than a single off-peak condition: Rail is organised through direction, station function and persistence around discrete node roles, while Bus follows a more continuous area-level activity--persistence gradient. These usage types are systematically, though not completely, associated with night-work geography, with correspondence strongest at the central and peripheral ends of both typologies. Intermediate and late-persistent types, most notably Rail C3, show temporal and network differences that are not fully captured by the area-level night-work classification and are associated with wider socio-spatial context.

This research makes three contributions. Empirically, it shows that a dedicated night-time window reveals structured temporal and spatial variation that is difficult to distinguish within all-day analysis. Comparatively, it demonstrates that Rail and Bus capture different dimensions of night-time mobility and require mode-sensitive rather than pooled interpretation. Methodologically, it links transport-derived temporal typologies with independently measured socio-spatial context, offering a reproducible framework for comparing mobility-based classifications with urban environmental data without one defining the other.

Practically, the resulting typologies provide a more differentiated basis than total ridership or centrality for identifying settings where questions of service continuity, interchange and time-dependent accessibility warrant closer attention. These include persistent but not necessarily central Rail nodes and outer areas where early decline coexists with resident employment in night-related sectors. Whether these settings reflect service constraints, mode substitution or lower travel needs remains a question for further research combining these typologies with service-supply, origin--destination and passenger-level data.
